\documentclass[11pt,a4paper]{article}
\usepackage[margin=22mm]{geometry}
\usepackage{amsmath,amssymb,amsthm,booktabs,longtable,array,fontspec}
\IfFontExistsTF{Noto Sans CJK SC}{\setmainfont{Noto Sans CJK SC}[AutoFakeBold=2,AutoFakeSlant=0.15]}{\setmainfont{Microsoft YaHei}[AutoFakeBold=2,AutoFakeSlant=0.15]}
\XeTeXlinebreaklocale "zh"
\XeTeXlinebreakskip=0pt plus 1pt
\setmonofont{lmmono10-regular.otf}[Scale=0.9]
\usepackage[colorlinks=true,linkcolor=blue,urlcolor=blue]{hyperref}
\usepackage{xurl,fancyvrb}
\DefineVerbatimEnvironment{verbatim}{Verbatim}{fontsize=\small}
\newcommand{\FS}{\operatorname{FS}}
\newcommand{\Z}{\mathbb Z}
\newcommand{\T}{\mathcal T}
\newcommand{\Def}{\mathcal D}
\theoremstyle{definition}
\newtheorem{theorem}{定理}[section]
\newtheorem{lemma}[theorem]{引理}
\renewcommand{\proofname}{证明}
\renewcommand{\abstractname}{摘要}
\renewcommand{\refname}{参考文献}
\renewcommand{\arraystretch}{1.16}
\setlength{\emergencystretch}{2em}
\linespread{1.08}
\raggedbottom
\title{将独立核查上界推进至 87\\\large 88 元端点的格化、共同稠密核心与完整证书接入}
\author{Sum-of-Subset 研究项目}
\date{2026 年 10 月 19 日}
\begin{document}
\maketitle
\begin{abstract}
上一轮独立证明链已给出 $71\le N\le88$。本报告核对归档中大小 88 的
结构归约及其全部必要模型行族，将该端点接到新的独立链之后。
证明方法的重点是把多个平移重叠统一放入同一实际缺额，使用一个同时适用的稠密子核心，
并且每个重叠组只计一次所选成员的外部损失。
有限部分复用冻结归档的整数证书，对全部普通根、九类补充证明及其精确补集重新运行。
归档材料本身已声称更强的 $N\le84$；本文推进的是独立核查链，不声称 87 优于归档数值 84，
也不声称重新发现了这些归档证书。一般证明的人工核对与有限程序的精确重放分别列出。
\end{abstract}

\section{结论与依赖边界}
设有限实数集 $A$ 满足 $\FS(A)\subseteq4A$，其中子集和包括空集，四重和允许重复。
以 $N$ 表示最大可能基数。
\begin{theorem}\label{main}
在此前已经独立核查的基础引理与 $N\le88$ 证明之上，大小 88 的端点被排除，故
\[\boxed{71\le N\le87.}\]
\end{theorem}
\begin{proof}
若一个 88 元可行集不含零，加入零仍可行，并得到与已证上界 88 冲突的 89 元集。
因此它必含零。沿用保留零坐标的整数化：保留选定子集和等式、元素互异及顺序。
第 \ref{lattice} 节在使用锐化剥离常数之前证明共同格和 $D$ 的整性。
第 \ref{rows}--\ref{cover} 节将每个实际配置置于一个必要模型及一个穷尽子分支。
第 \ref{verification} 节的完整精确证书排除全部这些模型。
下界 71 沿用此前已检验构造。
\end{proof}
这里不使用任意删去元素后的可行性，也不重新调用归档中更早的大小 89 排除结果。
新的降阶前提是提交 \texttt{a1ee834} 的独立 $N\le88$ 证明。

冻结证据位于仓库 \path{archive/upper86/} 的 28 个 ZIP 分片。
本文只接入其中 \path{proofs/exact87/} 的大小 88 端点，
其名字 87 指所得上界。归档一般引理的已知剥离勘误和中央区间端点勘误继续有效：
中央区间为 $[h+1,L-h-1]$，衔接条件为 $x\le L-2h-1$。

\section{先格化，再使用锐化预算}\label{lattice}
记两层剥离后
\[
 |D|=80,\quad |E|=72,\quad E=P\cup\{0\}\cup(-Q),\quad
 |P|=a,\ |Q|=b,\ a+b=71.
\]
令 $T_r=r(r+1)/2$。先取 $a\le b$，写成 $a=35-k,b=36+k$。
\[
 I=k(k+1)=(a-35)(a-36),\quad
 \Lambda=1+T_a+T_b=1297+I,
\]
\[
 \Omega=22\cdot80+97-\Lambda=560-I,\qquad
 C=2\cdot71+4\cdot72+14-\Lambda=-853-I.
\]
$\mu_X=|\FS(X)|-1-T_{|X|}$；$\rho_X$ 为此前定义的四极值缺陷，
满足 $\rho_X\le\mu_X$。记支持的洞数为 $h_P,h_Q$，和集超额为 $\kappa$，
$z=h_P+h_Q-\kappa$。沿用精确预算
\begin{equation}\label{budget}
 \rho_P+\rho_Q+2\vartheta+c+\T=\Omega,
 \qquad c=2|2D|+18\cdot80+85-|\FS(D)|.
\end{equation}
在共同格上证明 $D$ 整数化以后，才得到
\begin{equation}\label{post}
 w\le4z,\quad \T\ge\kappa+150-3z,\quad
 \T\ge\max\{50,z,150-2z\}.
\end{equation}
原来最外层的八个元素此时尚未自动整数化。

\subsection{边界与等差支持}
缺号侧的矛盾 $1+T_{72}=2629>2189$ 保证完整的八元剥离。
未配对数 $q$ 满足 $q^2\le2801$ 且为奇数，故 $q\le51$，先有 $a,b\ge10$。
大侧缺陷非负、小侧缺陷至少 $-8$，将边界缩到 $a\ge12$。
\begin{center}\begin{tabular}{rrrrr}\toprule
$a$&$b$&$\Omega$&$\kappa$ 的上限&修正后的内层比较上限\\\midrule
12&59&8&16&$-1349$\\
13&58&54&62&$-1165$\\
14&57&98&98&$-1005$\\\bottomrule
\end{tabular}\end{center}
各行先满足严格逆定理的两侧条件。删去最大元的前缀可能只有 11、12、13 元，
故用通用 $-8$ 下界，在内层比较中保留 $+16$ 修正；表中比较仍严格为负。
外层比较也严格成立。此后才能用 $\T\ge50$。
$a=12$ 立即排除；$a=13,14$ 若大侧非等差，用 $\rho_Q\ge b-5$ 排除；
若大侧等差，则共同格上的零洞强迫 $z=0$，用 $\T\ge150$ 排除。
最终 $15\le a,b\le56$。

等差支持的大侧至少有 $667$ 点，超过 $\kappa\le560$；剩余类及间隙公式给出共同格和 $z=0$。
若仅小侧等差，$\Omega<2T_a$ 且 $\Omega-T_a\le128$，最多两个剩余类。
单类长间隙造成最大元平移分离，缺陷矛盾的最小余量为 79。
双类情形先规范到半步格，以最小奇系数 $w_0$ 计质量。
$w_0>u=|\FS(P)|$ 的质量余量至少 1003；$w_0\ge15$ 时
$7b-289/8>2b+125$，故 $w_0\le13$，从而 $z\le12$。
使用该情形专属跨度估计，内层比较至多 $-955$，外层比较亦严格，遂证明 $D$ 整性。
现在 $\T\ge T_a+114$；还必须保留大侧非等差缺陷 $b-5$，得到
\[
 T_a+114+(b-5)-\Omega\ge28>0.
\]
单用 $T_a+114$ 不足以排除本端点。所有留下的等差情形均有 $z=0$。

\subsection{十二个不足的逆定理方向}
采用归档所需的较弱下界
$\nu(r)=r-5$（$r\ge21$）及 $r-14$（$15\le r\le20$）。
小侧充分条件
\[
 a-3+\nu(a)+\nu(b)>\Omega-T_a+2
\]
恰可能在 $k=7,\ldots,18$ 失败。大侧严格逆条件始终成立。
若小侧条件实际失败，令
\[
 R=\Omega-T_a+2-(a-3)-\nu(a),\quad
 K_{\max}=\Omega-\nu(a)-\nu(b),\quad U_{\max}=K_{\max}+3,
\]
\[
 U_{\min}=T_a+a-2,\qquad V_{\min}=T_b+b-2.
\]
先在 $Q$ 的内禀原始整数格上规范化，此时不预设 $P$ 整数化。
使用已独立核对的原始集合有效洞引理；给定原始性后，该引理不需要先假设洞数小于支持大小。
具体见证如下：
\begin{center}\small\begin{tabular}{rrrrrrrrrrrrr}\toprule
$k$&7&8&9&10&11&12&13&14&15&16&17&18\\
$\beta$&3&3&4&4&4&4&4&4&4&3&2&2\\
$C_h$&51&70&82&93&97&100&98&91&92&79&61&38\\
$R$&52&66&77&85&90&92&91&87&89&79&66&50\\\bottomrule
\end{tabular}\end{center}
每个见证都检验有效洞引理的两条严格参数不等式；若 $h_Q>2\beta$，
则 $\mu_Q\ge b+C_h+1$，旧的灵活缺陷逆界给出 $\rho_Q>R$。
因此 $h_Q\le H=2\beta\le8$。

若 $\FS(P)$ 占 $Q$ 格的 $t$ 个剩余类，则
$\kappa\ge(t-1)(|\FS(Q)|-1)$；$K_{\max}<V_{\min}-1$ 强迫 $t=1$。
令两侧支持直径为 $A_0,B_0$。$V_{\min}>U_{\max}+K_{\max}-1$ 排除 $A_0>B_0$；
两端平移的区间容器给出 $h_P\le\kappa+h_Q$。
接着短直径逆定理给出 $h_P,h_Q\le\kappa$、$0\le z\le H$。
该短直径定理不额外要求两侧分别原始。
四个严格余量的最小值依次为 $543,99,536,162$。

普通内层比较为 $-293-2I+2z$，故 $z\le146+I$ 已够。
若 $z\ge147+I$，预算先排除 $k\ge13$；其余两侧至少 23 元。
任一最终平移分离都会给出至少 30 的预算矛盾余量。
因此最大元删除后的洞至少保留一半，改用一次总预算得到比较上限 $-293-2I<0$。
外层比较为 $1297>4\cdot80-10+560$。这证明 $D$ 整性，随后才能调用式 \eqref{post}。
修复情形最后有 $\kappa\le T_a-36$，保证小侧原始性及后续严格洞假设；
严格逆情形已经具备它，等差情形 $z=0$。

\section{把全部覆盖行归入同一实际缺额}\label{rows}
令 $F_s$ 为正侧或反射后的负侧实际目标，$F_s\subseteq[1,S_s]_{\Z}$。
对一组带标签的分量 $A_i$，定义
\[
 \Def_F(J)=\sum_{i\in J}|A_i|-\left|\left(\bigcup_{i\in J}A_i\right)\cap F\right|.
\]
所有组变量均取这个实际值；单分量端点损失取实际下尾、上尾数。
不同辅助变量可以共享几何点，但在同一覆盖预算中不能重复使用容量标签。

\begin{lemma}[统一的缺额规则]\label{deficit}
若标签组 $J_1,J_2$ 不交，则
$\Def_F(J_1\cup J_2)\ge\Def_F(J_1)+\Def_F(J_2)$。
若 $B_i\subseteq A_i$，则对任一 $j\in J$，
\[
 \Def_F(J)\ge\sum_{i\in J}|B_i|-\left|\bigcup_{i\in J}B_i\right|+|A_j\setminus F|.
\]
若在标签上选取任意森林 $\mathcal G$，则
\[
 \Def_F(J)\ge\sum_{ij\in E(\mathcal G)}|A_i\cap A_j|+|A_j\setminus F|
\]
对任一单独所选成员 $j$ 成立。
\end{lemma}
\begin{proof}
第一式由并集次可加性直接得到。后两式逐点计数。
若一点属于 $m$ 个分量，在目标内贡献缺额 $m-1$，在目标外贡献 $m$。
子集合的原始重叠最多占 $m-1$，森林诱导边数也最多为 $m-1$；
目标外再加至多一次所选成员指示仍不超过 $m$。
零重数的点各项均为零。
\end{proof}
此规则同时解释单核心重叠、按顺序累积的核心组、前导组和稠密星形组，
并明确禁止把两条替代性下界相加。

\subsection{共同稠密子核心与星形计数}
记 $R_p=\max P,R_q=\max Q$，四个剥离正元为
$R_p<x_1<x_2<x_3<x_4=M_p$，负幅值相应为 $y_i$。
定义
\[
 g_p=\max\{0,\max_{1\le j\le4}(2j+R_p-x_j)\},\quad
 g_q=\max\{0,\max_{1\le k\le4}(2k+R_q-y_k)\}.
\]
整数间隙保证 $0\le g_p,g_q\le4$。分别选取达到最大值的 $j,k$，允许取零，
将两侧对应前缀同时加入 $E$，得到\emph{同一个} $E'\subseteq D$，满足
\[
 2|E'|-\operatorname{diam}E'=2e-R_p-R_q+g_p+g_q,\qquad e=72.
\]
正侧前导分量 $A_0=c_p+2D$，其中 $c_p=p_1-p_2-p_3-p_4$。
对保留的自身分量 $A_i=c_p+p_i+D$，比较前导中的 $c_p+M_p+E'$ 与叶中的 $c_p+p_i+E'$，得到
\[
 |A_0\cap A_i|\ge143-R_p-R_q-p_i+M_p+g_p+g_q.
\]
交叉分量 $c_p-q_j+D$ 使用前导中的 $c_p-M_q+E'$，将最后两项换为 $-q_j+M_q$。
所需的是相对位移整性。前导中心的绝对位置不影响原始交集。
反射给出另一侧。

令各交集辅助变量取真实交集数，用引理 \ref{deficit} 的星形森林得到
\[
 \Def_F(\{0,1,\ldots,t\})\ge
 \sum_{i=1}^t|A_0\cap A_i|+|A_j\setminus F|\quad(0\le j\le t).
\]
不同交集可相互重叠，无需假设它们不交。每个 $j$ 给同一缺额的替代下界。
前一轮只利用一个自身叶的核心重叠；此处一个同时实现的稠密核心可加强所有叶，并允许全部叶一起计数。

\subsection{逐行接入的门槛与实际赋值}
冻结模型的所有一般行族按以下清单核对；完整路径和哈希另存于接入清单。
\begin{longtable}{p{.26\linewidth}p{.66\linewidth}}
\toprule 行族&必要性与接入条件\\\midrule\endhead
基本顺序、质量、跨度&$D$ 的整数间隙；$\sum P=T_a+\mu_P+h_P$；四个剥离点的最小和；最大元删除后的旧缺陷界。正洞质量保留同一个 $z$。\\
十八单元与单分量损失&按实际贡献者定向；配对单元最多记一次容量。A 分支的特殊 beta 配对保留两种实际贡献者变体。\\
全局、核心、普通分位界&常数分别为低/高端 $80/79$、$76/75$、正侧 $b+5/a+4$。核心门槛严格为残余损失 $<76$，普通门槛严格为 $<\min(a,b)+5$。\\
双和集包络&内部间隙辅助量取间隙与区间宽度的最小值；二尾计数、混合尾及固定删除包络沿用已证容器计数，删除参数重新代入 $d=80$。\\
自身配对与顺序核心组&不同匹配边使用不交标签；每块用引理 \ref{deficit} 的子核心容器，保留一个成员损失；前缀组缺额作一致实际赋值。\\
前导核心组&使用 $2E$ 的至少 $2e-1$ 点与若干 $E$ 叶；原始容器计数后只保留一个损失。\\
完整 $D$ 组与截断&前导容量 $K=|2D|$ 只记一次。原始区间容器可加一个损失，截断到目标的容器不再额外加该损失。所选组包含实际贡献的自身叶，故截断长度上下界合法。\\
极端核心与星形行&原始交集需相对整性；两个符号的组使用不同容量标签。小的标量预算行可丢弃其余非负损失。\\
固定分数层&110 个不同多项式，26 个不交非正组。额外高端标签排除已计入的强制损失标签。仅在该根的整数左逆通过后使用整数分支。\\
补充量化、洞区间&固定损失只收在覆盖预算中；正洞缺陷改善 18；三个洞区间使用严格有效洞见证。模型旧行及原始格性门槛保留。\\\bottomrule
\end{longtable}
特别地，基础正洞模型使用较弱 $\rho_P+\rho_Q\ge71-28=43$。
当两侧至少 21 时可用 $61$，额外改善 18。不能在不满足条件的根上全局替换。

\section{格性、穷尽分类与补充分支}\label{cover}
\subsection{不能混淆三种格性结论}
三空单元时残余损失可能为 170，大于 $|2D|\ge159$，故不能用“前导必贡献”的捷径。
全部 2992 种模式由实际保留中心矩阵 $C$ 得到整数 $W$，逐项验证 $WC=I_8$。
于是 $Cx$ 整数推出八个外元 $x$ 整数，与前导是否贡献无关。
至多两空的解析论证是其前置情况。

四空、五空时相应上限足够小，两前导均贡献，从而相对位移整数；
这本身仍不推出缺失外元整数。需要严格外元间隙或固定分数二分的补充根，
另有该根自己的 $WC=I_8$ 证书。所有左逆均来自原模型及原门槛，
后加的损失改进不能倒过来改变格性来源。
前导、嵌套前导和核心分位界的门槛分别严格为 $159,151,76$。

\subsection{完整根域}
令 $B_0=410-I-80j$。低洞域为 $Z:z=0,j=0,\ldots,5$，
和 $P:1\le z\le14,j=0,\ldots,4$；后者的保守上限为 $B_0-15$。
共同必要行包括
\begin{align*}
 \delta&\le B_0\quad(Z),&\delta-2z&\le B_0-43\quad(P),\\
 R_p+R_q+\delta&\le75+B_0+4z,&
 S_p+S_q+\delta-2K&\le1524-80j+4z.
\end{align*}
这里 $\delta=c-80j$；高洞仍保留同一个 $z$，不得分别最大化洞数和预算后拼接。
\begin{center}\small\begin{tabular}{lrrrrrr}\toprule
$j$&0&1&2&3&4&5\\\midrule
模式数&4&70&578&2992&10880&29512\\
$Z$ 根&160&2520&18496&77792&217600&177072\\
$P$ 根&160&2520&17340&71808&195840&0\\\bottomrule
\end{tabular}\end{center}
低洞共 781308 根。高洞枚举 $0\le k\le20$、$15\le z\le560-k(k+1)$，
共 8386 个原始标量状态。630 个有效洞参数逐项检验严格条件。
剔除已被全活跃质量矛盾排除的状态后，保留 9327 个状态/空数对，合并为 188 个闭整数区间。
各区间的支持下界和 $\max(50,z,150-2z)$ 分段保持不变。
两反射方向均列出，因为 $a+b=71$ 为奇数。高洞共 346972 根，总数 1128280。

\subsection{十四单元损失覆盖与正洞分段}
八个零洞父根为 A 分支、缺失两配对及两最大自身单元，$32\le a\le39$。
每侧保留七个实际非配对平移。缺失最大元论证给出八个额外固定高端标签，
扣除后残余上限依次为
\[
 70,76,80,82,82,80,76,70.
\]
若 $C<76$，全部十四个整数损失 $L_i\le75$。
否则完整覆盖由“全部 $\le75$”及十四种“指定 $L_i\ge76$”组成；
$C<152$ 保证至多一个大损失。
指定大损失时，其余损失至多 $C-76<\min(a,b)+5$，故只有其余单元可追加普通分位界。
两端各需一子模型，其余六根各需十五个，共 92 个子模型。
一般贡献行 $v+M_s\ge1$、$v-M_t\le S_s$、$L_v\le79$ 只用于实际非配对贡献者。

十四个正洞四空父根则为 $29\le a\le42$，相同缺失模式。
每个父根都必须覆盖三个闭整数区间 $[1,2],[3,4],[5,14]$。
对区间下端 $l$，仅当 $l>2\beta$ 才调用有效洞见证。
28 个 $(r,\beta)$ 见证和 42 个支持下界输入均由程序检查。
改变质量行时保持其洞系数，保留旧预算标签及格性来源。
此处正洞的强制缺失最大元固定损失为零，不能借用零洞的 8。

\subsection{所有二分的数学来源}
普通二分只分“普通分位界成立”或“位置越界且有相应尾损失”；
内层二分只分“至少 $j$ 个尾点”或“第 $j$ 个内层点在目标另一侧”；
核心二分只分“两个核心分位界成立”或“总损失至少 76”。
特殊配对要么另一成员实际贡献，要么所选成员是唯一贡献者；只有后一分支按原生容量加入所选成员损失。
固定分数用整性的 $f\le S$ 或 $f\ge S+1$、$f\le0$ 或 $f\ge1$。
所有辅助缺额是实变量；不对未经证明为整数的辅助量作整数二分。
每个使用的二分都保留两个孩子。

\section{精确证书、复现和实际运行}\label{verification}
普通证书池有 10689 个规范正整数向量。
固定顺序的 1128280 个无符号 16 位项共 2256560 字节；631 个零项恰是补充父根。
其余项选择向量和 fast/full 模型。每次从冻结源码重新生成当前矩阵，检查
\[
 y\ge0,\quad y^{\mathsf T}A\ge0,\quad y^{\mathsf T}b<0,
 \qquad Ax\le b,\ x\ge0.
\]
对照的是完整根描述的哈希及实际父根集合，不只检查计数相同。
\begin{center}\small\begin{tabular}{lrrrr}\toprule
补充类别&父根&子模型&二分&叶\\\midrule
零洞三空&20&20&24&44\\
零洞四/五空&68&68&240&308\\
零洞损失覆盖&8&92&2510&2602\\
直接正洞缺陷&106&106&0&106\\
正洞三空&26&26&59&85\\
正洞三段覆盖&14&42&417&459\\
高洞第一组&41&41&91&132\\
高洞第二组&170&170&376&546\\
高洞第三组&178&178&635&813\\\midrule
合计&631&743&4352&5095\\\bottomrule
\end{tabular}\end{center}
普通精确反证 1127649 个，加上补充叶 5095 个，共 1132744 个；
恒等式为 $1132744=1128280+4352+112$，其中 112 是两类有限覆盖额外增加的子模型数。

本次新增入口\path{verification/upper87.py} 校验原归档分片及拼接哈希，
仅提取大小 88 端点；执行全部结构、分类、九类补充及精确补集检查，
另将普通根范围完整分片重放。分片区间首尾相接、无重无漏，且源与根清单身份一致才接受。
这不是一次默认串行调用，但覆盖与它相同的全部有限义务。
\begin{verbatim}
python -I -S -B verification/upper87.py --jobs 8 \
  --work-dir .work/reproduce87 --output result87.json
python -I -S -B verification/upper87_geometry.py \
  --output geometry87.json
\end{verbatim}
Windows 可用\path{verify_upper87_windows.bat}，或将上述 python 换为 \texttt{py -3}。
新入口拒绝 \texttt{-O}，因为冻结归档的部分检查使用断言；这次拒绝已实际测试。
几何程序不用断言，标准与 \texttt{-O} 输出逐字节一致。
复现时须使用新的工作目录，旧成功日志不是输入。

\paragraph{本次运行回执。} 2026 年 10 月 19 日，8 个普通范围与完整补充阶段全部实际退出为 0；最后一个范围完成后重新检验冻结清单。总体状态为\texttt{PASS\_COMPLETE\_ENDPOINT\_FINITE\_REPLAY}，总用时 556.933 秒（本环境并发运行，不作为跨机器性能承诺）。全部 1128280 个原始根覆盖完整，最终 1132744 个叶反证通过。来源清单中的 38 个接入文件哈希与本次实际执行副本一致。另有 11 项算术/入口负面控制、缺失范围及重复范围两项拼合负面控制通过。

\noindent\begin{minipage}{\linewidth}
原归档 SHA-256 为
\begin{center}\scriptsize\texttt{390d390c1172471ff97f04113d54937de2ae070cabcb272f5529f08a19a069df}\end{center}
端点清单 SHA-256 为
\begin{center}\scriptsize\texttt{9e3631ce91c9ab07caba32b8ed55c3b0c011870ebf6a5b887206ecd3cab86c75}\end{center}
\end{minipage}\par\medskip

逐范围命令、真实退出码、各阶段原始日志、源码身份与输出哈希均随结果包保存。


新增直接集合检查使用固定种子 2026101901：240 个实际有限集合中，
150 个出现严格稠密增益；检查 240 个共同稠密核心恒等式、2880 个原始交集界、
12960 个星形行、480 个资源分配行。
此外对最多八叶星形的全部逐点成员关系检查 16384 次，
并明确拒绝“加两个外部损失”的反例。
这些有限例子验证实现，不替代第 \ref{rows} 节对任意集合的证明。

\section{来源与本次接入范围}
本次人工核对涵盖：大小 88 的边界、等差/半步格、十二个逆定理修复、
两层格化、低高洞分类、所有基本和几何行族、九类补充所用加强行、
实际贡献者和整数分支的许可条件。冻结纯函数与归档数学函数的 AST 一致性、
当前矩阵和所选叶的对应由完整程序再次检查。
这构成在既有独立基础引理上的数学接入，而非单由保存的 PASS 日志推断定理；
它仍不是证明助手中的形式化证明，也未完成更小的 87、86、85 元端点接入。

\begin{thebibliography}{9}
\bibitem{prior} 本仓库，\path{reports/upper88_leading_overlap_20261018.tex}，
独立 $N\le88$；以及其结构前提\path{reports/upper88_progress_20261018.tex}。
\bibitem{audit} 本仓库，\path{reports/proof_audit_20261018.md}，基础引理和归档证据边界。
\bibitem{archive} 冻结归档，\path{proof/complete_upper86_proof.tex} 的
\texttt{lattice88} 与 \texttt{upper87} 两节；
\path{payload/subset_sum_verification/proofs/exact87/} 中的必要模型、补充证明与纯算术验证器。
\bibitem{adoption} 本次机器可读来源清单，\path{data/upper87_adoption_sources.json}；
运行证据位于\path{results/2026-10-19/upper87/}。
\end{thebibliography}
\end{document}
